On the other hand, the relation
\(\textup{(1)}\times\textup{(1')}=
\textup{(2)}\times\textup{(2')}\) gives at once
\[
  aa'(d-1)^2=bb'(d+1)^2.
\]
Combining the last two relations immediately yields the second
formula of \SGARefLink{sga3:xix:lemma:5:3}{Lemma 5.3}.

\par\medskip\noindent
\textbf{5.4.}---
Now consider
\[
  G=G_1\cap\mathrm{SL}(\mathfrak g).
\]
This is the closed subgroup of \(G_1\) defined by the equation
\(\det(u)=1\); hence it is an affine group over \(S\).
\SGARefTarget{sga3:xix:paragraph:5:4}

\par\medskip\noindent
\SGAInlineRefTarget{sga3:xix:proposition:5-5}\textbf{\textit{\SGARefLink{sga3:xix:proposition:5:5}{Proposition 5.5}.}}---
\textit{The group \(G\) is smooth over \(S\).}
\SGARefTarget{sga3:xix:proposition:5:5}

To prove the proposition, we shall need the following lemmas.

\par\medskip\noindent
\textbf{\textit{\SGARefLink{sga3:xix:lemma:5:6}{Lemma 5.6}.}}---
\textit{Let \(U\) be the open subset of}
\[
  \operatorname{End}_A(\mathfrak g)
  \simeq
  \mathbb W\bigl(M_3(\mathcal O_S)\bigr)
\]
\textit{defined by the condition that \(ad\) be invertible; that is,
the principal open subset
\(\operatorname{End}_A(\mathfrak g)_f\), where
\(f(u)=ad\).  Then \(U\cap G\) is the closed subscheme of \(U\)
defined by the six equations}
\[
  \textup{(1)},\quad
  \textup{(2)},\quad
  \textup{(2')},\quad
  \textup{(3)},\quad
  \textup{(3')},\quad
  \textup{(D)}\colon aa'-bb'=d.
\]
\SGARefTarget{sga3:xix:lemma:5:6}

It is clear that these six relations are satisfied by every point
of \(G\), by \SGARefLink{sga3:xix:lemma:5-3}{Lemma 5.3}, and therefore by every point of \(U\cap G\).
Conversely, suppose that \(u\in U(S')\), for some \(S'\to S\),
satisfies the six conditions in the statement.  We must show that
\(\det(u)=1\), and that \(u\) also satisfies \textup{(1')},
\textup{(4)}, \textup{(4')}, and \textup{(5)}.

First, \textup{(D)} implies \textup{(5)}.  Relations \textup{(2)}
and \textup{(2')} give
\[
  bb'(d+1)=bce'=b'c'e.
\]
On the other hand, from \textup{(3)} and \textup{(3')}, writing
\(2t(bc-ac')(b'c'-a'c)\) in two ways, we obtain
\[
  (bc-ac')e'=(b'c'-a'c)e.
\]
Together with the preceding relation, this gives
\[
  ac'e'=a'ce.
\]
By \textup{(1)}, however,
\[
  a'ce=a'a(d-1).
\]
Thus
\[
  a\bigl(a'(d-1)-e'c'\bigr)=0,
\]
and this implies \textup{(1')}, since \(a\) is invertible.

Consequently, \textup{(1)}, \textup{(2)}, \textup{(2')}, and
\textup{(1')} hold.  \SGARefLink{sga3:xix:lemma:5:3}{Lemma 5.3} and \textup{(D)} now give
\[
  d\bigl(\det(u)-1\bigr)=0.
\]
Since \(d\) is invertible, it follows that \(\det(u)=1\).

It remains to prove \textup{(4)} and \textup{(4')}.  We prove
\textup{(4')}; the other calculation follows by symmetry.  From
\textup{(3)}, \textup{(3')}, and \textup{(D)}, we have
\[
  a'e+be'=-2t(aa'-bb')c'=-2tdc'.
\]
Combining \textup{(1')} and \textup{(2)}, we also have
\[
  a'e+be'=-d(be'-ea'),
\]
\footnote{Editor's note: the sign has been corrected.}
which proves \textup{(4')}, since \(d\) is invertible.

\par\medskip\noindent
\SGAInlineRefTarget{sga3:xix:lemma:5-7}\textbf{\textit{\SGARefLink{sga3:xix:lemma:5:7}{Lemma 5.7}.}}---
\textit{The group \(G\) is smooth over \(S\) along the identity
section.}
\SGARefTarget{sga3:xix:lemma:5:7}

